% Generated draft: requires mathematical and source review.
\documentclass[11pt]{article}
\usepackage[margin=25mm]{geometry}
\usepackage{fontspec}
\setmainfont{Latin Modern Roman}
\usepackage{amsmath,amssymb,mathtools,unicode-math}
\setmathfont{Latin Modern Math}
\usepackage{enumitem,needspace,xurl,hyperref}
\setlength{\parindent}{0pt}
\setlength{\parskip}{5pt}
\setlength{\emergencystretch}{3em}
\newcommand{\N}{\mathbb{N}}
\newcommand{\Z}{\mathbb{Z}}
\newcommand{\Q}{\mathbb{Q}}
\newcommand{\R}{\mathbb{R}}
\newcommand{\C}{\mathbb{C}}
\newcommand{\F}{\mathbb{F}}
\newcommand{\A}{\mathbb{A}}
\newcommand{\PP}{\mathbb{P}}
\newcommand{\E}{\mathbb{E}}
\newcommand{\eps}{\varepsilon}
\newcommand{\dd}{\,\mathrm{d}}
\begin{document}
\section*{UnsolvedMath ID: 2302040}
\textbf{Source problem code:} catalogue entry
\subsection{Definitions and mathematical statement}
For a nonconstant entire function $f$, define
\[
M(r,f)=\max_{|z|=r}|f(z)|,\qquad
E(c)=\{z\in\mathbb C:|f(z)|>c\},
\]
and let $m_2$ be planar Lebesgue measure. Assume $m_2(E(c))<\infty$ for some real threshold $c$.

\textbf{Growth question.} Determine the minimum possible growth of $f$ under this hypothesis. Hayman's proposed necessary and best-possible condition is
\[
\int_R^{\infty}\frac{r\,dr}{\log\log M(r,f)}<\infty,
\]
where $R$ is sufficiently large that $M(r,f)>e$ for all $r\geq R$. This tail-integral formulation makes sense of the source's formal lower limit $0$, since $\log\log M(r,f)$ need not be positive or defined near the origin. Finiteness does not depend on which sufficiently large $R$ is used. The phrase best possible is retained, but the source gives no explicit converse construction or admissible growth-envelope class defining its sharpness.

\textbf{Threshold question (P. Erd\H{o}s).} If $m_2(E(c))<\infty$, must $m_2(E(c'))<\infty$ also hold for smaller thresholds $c'<c$? The substantive positive-threshold version asks this for every $0<c'<c$. The source does not explicitly state that restriction; nonpositive thresholds introduce a separate degeneracy because a nonconstant entire function has only a discrete zero set and $E(c')$ then has infinite area. This missing qualification is not silently ignored.

The supplied summary labels the problem open as of Hayman's 2018 edition and reports no recent resolution located, without an independently verified current result.
\subsection{Short English statement}
How fast must an entire function grow if the region where its modulus exceeds one threshold has finite area? Test the proposed integral condition and ask whether finite area persists at every smaller positive threshold.
\subsection{Sources}
\begin{itemize}
\item Hayman, \emph{Research Problems in Function Theory} (2018), Problem 2.40, attributed to P. Erd\H{o}s, with the growth conjecture explicitly attributed to Hayman, \url{https://arxiv.org/abs/1809.07200}; source-TeX extraction, accessed 2026-07-29. The 2018 update was retained for status triage.
\item Local registry \texttt{data/unsolvedmath/problems.json}, permanent ID $2302040$, source code catalogue entry, \emph{Research Problems in Function Theory --- Problem 2.40}. The record labels this open and OPEN-TRIAGE, with NEEDS\_REVIEW source-status and rights qualifications. No separate sharpness or threshold-result citation is supplied.
\end{itemize}
\end{document}
